% !PRINTED 258
\weberpage{258}
%<<CONT

\[
f(\omega) + \frac{\sqrt{2}}{f(\omega)}
\]

abhängen. Sie sind \gesperrt{ganze} Funktionen dieser Verbindungen,
wenn die $\Phi_{a,c,\partial}$ für jedes endliche, nicht reelle $\omega$ endlich bleiben,
und sie sind \gesperrt{konstant}, wenn alle $\Phi_{a,c,\partial}$ auch für $q = 0$, d.\,h.
für $f(\omega) = 0$ und $f(\omega) = \infty$, endlich bleiben. In diesem Falle
sind \gesperrt{sämtliche $\Phi_{a,c,\partial}$ einer und derselben Konstanten
gleich} (\S\,70).

Daraus ergibt sich der Satz:

\gesperrt{Bildet man ganze rationale Funktionen $\Phi_{a,c,\partial}$ aus}

\[
f\left(\frac{c+\partial\omega}{a}\right), \quad f(\omega), \quad \frac{1}{f\left(\dfrac{c+\partial\omega}{a}\right)}, \quad \frac{1}{f(\omega)},
\]

\gesperrt{welche die Eigenschaft haben:}

1. \gesperrt{durch die Substitutionen}

\[
\left(\omega, -\frac{1}{\omega}\right), \quad (\omega, \omega+2), \quad \left(\omega, \frac{\omega-1}{\omega+1}\right)
\]

\gesperrt{nur untereinander vertauscht zu werden,}

2. \gesperrt{für $q = 0$ nicht unendlich zu werden, so ist}

\[
\Phi_{a,c,\partial} = \text{constans}
\]

\gesperrt{eine Transformationsgleichung.}

Um diese Bedingungen zu befriedigen, ist zunächst der Einfluß
der Substitutionen (1) auf die Funktionen $f$ zu untersuchen.
Dieser ergibt sich aus den Zusammensetzungen \S\,69, (8) bis (13)
und aus den Transformationsformeln für die $f$-Funktionen \S\,40.

Zur Abkürzung führen wir die Bezeichnung ein:

\weq{3}{\begin{aligned}
f(\omega) &= u, & f\left(\frac{c+\partial\omega}{a}\right) &= v, \\[4pt]
f_1(\omega) &= u_1, & \left(\frac{2}{a}\right) f_1\left(\frac{c+\partial\omega}{a}\right) &= v_1, \\[4pt]
f_2(\omega) &= u_2, & \left(\frac{2}{\partial}\right) f_2\left(\frac{c+\partial\omega}{a}\right) &= v_2.
\end{aligned}}

Es ergeben sich dann folgende zusammengehörige Änderungen,
wenn in der Bezeichnung auf die Verwandlung der
Zahlen $a, c, \partial$ in $a, c_1, \partial$ oder $a_2, c_2, \partial_2$, wie sie eben durch die
angeführten Formeln des \S\,69 charakterisiert ist, keine Rücksicht
genommen wird:
%CONT>>
